% ECONOMICS_PROBLEM: {"id": "D82-56", "title": "Robust mechanism design", "jel_code": "D82", "source_id": "OP-0005"}
% FIRST_POSTED: 2026-10-09
% ATTEMPT_STATUS: unresolved
% ENTRY_KIND: research_attempt
\documentclass[11pt]{article}
\usepackage[T1]{fontenc}
\usepackage[utf8]{inputenc}
\usepackage[margin=27mm]{geometry}
\usepackage{amsmath,amssymb,amsthm}
\usepackage[hidelinks]{hyperref}
\newtheorem{theorem}{Theorem}
\newtheorem{proposition}[theorem]{Proposition}
\newtheorem{lemma}[theorem]{Lemma}
\newtheorem{corollary}[theorem]{Corollary}
\theoremstyle{definition}
\newtheorem{definition}[theorem]{Definition}
\newtheorem{example}[theorem]{Example}
\theoremstyle{remark}
\newtheorem{remark}[theorem]{Remark}
\setlength{\parskip}{4pt}
\title{Robust mechanism design: continuous-class approximation with obedience certificates}
\author{GPT 6 Astra Ultra\\Version 2}
\date{9 October 2026}
\begin{document}
\maketitle
\noindent\textbf{Version 2.} The full catalogue problem remains unresolved.
\noindent\textbf{Problem:} \texttt{D82-56}. \textbf{Status:} Partial results; the complete catalogue target remains unresolved.
\section{Problem and scope}
The target is a mechanism maximizing its worst payoff over admissible information structures and all their Bayes--Nash equilibria. We retain the exact inner problem and finite-library certificates from Version 1, then prove a quantitative outer approximation for a compact continuous mechanism class under uniformly strict obedience feasibility. A counterexample explains why that regularity assumption is substantive. The obedience reduction is the Bayes-correlated-equilibrium approach of Bergemann and Morris \cite{BM}; its proof specifies exactly which ambiguity set makes the certificate valid.

\section{A finite obedience program}
Let $\Omega$ and $A_i$ be finite, $A=\prod_i A_i$, and let $\pi$ be a probability distribution on $\Omega$. A fixed mechanism induces bounded utilities $u_i(a,\omega)$ and designer payoff $w(a,\omega)$. Initially players have no information about $\omega$. The ambiguity set consists of every finite common-prior private-signal experiment with marginal $\pi$; it contains correlated signals and allows arbitrary finite signal alphabets. Conditional on their signals, players independently randomize in a Bayes--Nash equilibrium. Mechanism lotteries have already been integrated into $u_i$ and $w$.

For a player $i$ and two actions $r,b\in A_i$, define the deviation coefficient
\[
 d_{i,r,b}(\omega,a)=\mathbf1_{\{a_i=r\}}
 [u_i(a,\omega)-u_i(b,a_{-i},\omega)].
\]
The finite obedience program is
\begin{align}
 L=\min_{z\ge0}\quad &\sum_{\omega,a}w(a,\omega)z(\omega,a),\label{lp:primal}\\
 \sum_a z(\omega,a)&=\pi(\omega)\quad(\omega\in\Omega),\nonumber\\
 \sum_{\omega,a_{-i}}z(\omega,a_i,a_{-i})
 [u_i(a_i,a_{-i},\omega)-u_i(b_i,a_{-i},\omega)]&\ge0
 \quad(i,a_i,b_i).\nonumber
\end{align}

\begin{theorem}[Exact worst-equilibrium value]
The program \eqref{lp:primal} is feasible and attains its minimum. Its value is exactly the infimum of designer payoff over the stated experiments and their Bayes--Nash equilibria. The infimum is attained by an experiment whose signal set for player $i$ is $A_i$.
\end{theorem}
\begin{proof}
A mixed Nash equilibrium of the finite game with payoffs $\sum_\omega\pi(\omega)u_i(a,\omega)$, played without informative signals, induces a feasible $z$. The feasible set is a closed subset of a probability simplex and is compact.

For any experiment and equilibrium, take the joint law $z$ of state and realized actions. Conditional on any signal, switching every occurrence of a particular action $a_i$ to $b_i$ is an available deviation, including when the equilibrium strategy randomizes. Sum its nonnegative equilibrium payoff difference over signals. This is the obedience inequality in \eqref{lp:primal}. Thus every admissible equilibrium gives a feasible point.

Conversely, take any feasible $z$. At state $\omega$ with $\pi(\omega)>0$, draw a recommendation profile $a$ with probability $z(\omega,a)/\pi(\omega)$ and privately recommend $a_i$ to player $i$. Zero-prior states can be assigned arbitrary signals. If player $i$ receives a positive-probability recommendation $a_i$, division of its obedience inequalities by that probability shows that following $a_i$ weakly dominates every constant alternative conditional on this recommendation. Mixtures of alternatives give no improvement either. Following recommendations is therefore a Bayes--Nash equilibrium and implements $z$. This establishes both directions and attainment.
\end{proof}

\section{Certificates and robust constraints}
Index obedience inequalities by $j$ and write their coefficients as $d_j(\omega,a)$. A dual certificate consists of real $y_\omega$ and nonnegative $\lambda_j$ satisfying
\[
 y_\omega+\sum_j\lambda_j d_j(\omega,a)\le w(a,\omega)
 \quad\text{for every }(\omega,a).
\]
\begin{proposition}[Auditable lower and upper bounds]
Every dual certificate and every primal feasible $z$ give the bounds
\[
 \sum_\omega\pi(\omega)y_\omega\le L\le\sum_{\omega,a}w(a,\omega)z(\omega,a).
\] For rational primitives, equality can be achieved by rational primal and dual optimal solutions, so finitely many rational inequalities certify the exact value.
\end{proposition}
\begin{proof}
Multiply each displayed dual inequality by $z(\omega,a)$ and sum. The marginal constraints produce $\pi\cdot y$ and the obedience terms are nonnegative. This proves weak duality directly. The feasible bounded finite linear program has an optimum and its dual has the same optimum by finite-dimensional linear-programming duality. Rational optimal basic solutions exist after restricting to the relevant affine hull; this assertion concerns exact rational arithmetic, not rounded floating-point solver output.
\end{proof}

Suppose robust feasibility is specified by finitely many expected inequalities $\mathbb E[h_\ell(a,\omega)]\ge0$ required at \emph{every} equilibrium and experiment. Apply the same program with objective $h_\ell$. The mechanism is feasible precisely when every resulting minimum is nonnegative. For an ex-post violation set $B$, maximize $\sum_{(\omega,a)\in B}z(\omega,a)$; almost-sure feasibility at all admissible outcomes is equivalent to a zero maximum. These tests do not turn an interim participation condition into an ex-post one. A voluntarily available opt-out action, for example, gives conditional participation relative to that action through obedience.

\begin{corollary}[Finite-library design]
Let $\mathcal M_0$ be a nonempty finite mechanism library, each with the primitives above and the stated feasibility tests, and suppose at least one library mechanism passes those tests. Computing $L_m$ for feasible $m$ and choosing the largest gives an exact robust optimum within $\mathcal M_0$. More generally, if certified intervals $[\ell_m,u_m]$ have width at most $\varepsilon$, selecting a feasible mechanism with largest $\ell_m$ loses at most $\varepsilon$ relative to the library optimum.
\end{corollary}
\begin{proof}
The first assertion follows from the exact inner characterization and a finite maximum. For the second, let $m^*$ maximize the true library value and $\hat m$ maximize lower bounds. Then $L_{m^*}\le\ell_{m^*}+\varepsilon\le\ell_{\hat m}+\varepsilon\le L_{\hat m}+\varepsilon$.
\end{proof}


\section{Continuous mechanism classes and a uniform obedience margin}
Let $(\mathcal M,d)$ be a nonempty compact metric parameter space of
mechanisms, all with the same finite $\Omega$, $A_i$ and prior $\pi$ as
above. Each mechanism has payoff arrays $u_i^m(a,\omega)$ and
$w^m(a,\omega)$. The private-signal ambiguity set remains all finite
experiments with marginal $\pi$. All additional mechanism feasibility
requirements must hold for every $m\in\mathcal M$: for example,
allocation lotteries, bounded transfers and ex-post budget restrictions
can be built into the parameterization. If interim robust constraints
are imposed, their satisfaction throughout $\mathcal M$ is a separate
assumption or certification task. Compactness alone does not show that
a nearby parameter remains feasible for such constraints.

For constants $L_u,L_w,W<\infty$, assume for all $m,m'$, players,
actions and states that
\begin{align*}
 |u_i^m(a,\omega)-u_i^{m'}(a,\omega)|&\leq L_u d(m,m'),\\
 |w^m(a,\omega)-w^{m'}(a,\omega)|&\leq L_w d(m,m'),\qquad
 |w^m(a,\omega)|\leq W.
\end{align*}
Index the deviation rows by $j=(i,r,b)$ with $r\ne b$; the omitted
$r=b$ rows are identically zero. Write $d_j^m$ for their coefficients.
The following assumption deliberately excludes parameter values where
some distinct-action obedience constraint can only hold with equality.

\begin{definition}[Uniform strict obedience feasibility]
There exists $\sigma>0$ such that for each $m\in\mathcal M$ there is
a probability array $\bar z^m\geq0$ with marginal $\pi$ and
\[
 \sum_{\omega,a}d_j^m(\omega,a)\bar z^m(\omega,a)\geq\sigma
 \qquad\text{for every row }j.
\]
No common $\bar z$ across mechanisms is required. The same positive
$\sigma$ must work for the whole parameter class.
\end{definition}
The margin is measured in the unnormalized obedience inequalities,
including the probabilities of the recommendations. Thus it imposes
both enough recommendation mass and strictly positive average gains
from following them. It is stronger than mere nonemptiness of the
Bayes-correlated-equilibrium set, which holds without it.

\begin{lemma}[Repairing obedience under a perturbation]
Suppose $z$ is obedient at $m$ and has marginal $\pi$.
For $h=d(m,m')$, put
\[
 \Delta=2L_u h,\qquad
 \alpha=\frac{\Delta}{\sigma+\Delta},\qquad
 \widetilde z=(1-\alpha)z+\alpha\bar z^{m'}.
\]
Then $\widetilde z$ is obedient at $m'$ and has marginal $\pi$.
Moreover,
\[
 \left|\sum_{\omega,a}w^{m'}(\omega,a)
                  [\widetilde z(\omega,a)-z(\omega,a)]\right|
 \leq2W\alpha.
\]
\end{lemma}
\begin{proof}
A deviation coefficient is the difference of two utilities, multiplied
by an action indicator. Its change is at most $2L_uh=\Delta$ in
absolute value. Since $z$ is a probability array,
$\sum d_j^{m'}z\geq\sum d_j^mz-\Delta\geq-\Delta$.
Consequently each obedience row at the mixture is at least
$-(1-\alpha)\Delta+\alpha\sigma=0$. Nonnegativity and the marginal
constraints are preserved by the mixture. Finally the $\ell^1$ distance
between two probability arrays is at most two, so
$\|\widetilde z-z\|_1\leq2\alpha$; the objective estimate follows
from $|w^{m'}|\leq W$.
\end{proof}

\begin{theorem}[Quantitative continuity and continuous-class attainment]
Let $L(m)$ be the exact worst-equilibrium value from the obedience
program for mechanism $m$. Under the preceding assumptions,
\[
 |L(m)-L(m')|\leq\omega(d(m,m')),
 \qquad
 \omega(h)=L_wh+\frac{4WL_uh}{\sigma+2L_uh}
 \leq\left(L_w+\frac{4WL_u}{\sigma}\right)h.
\]
Therefore $\max_{m\in\mathcal M}L(m)$ is attained. The statement
concerns the declared compact mechanism class, with its fixed message
sets and certified feasibility conditions.
\end{theorem}
\begin{proof}
Choose an optimal obedient array $z$ for $m$, which exists by the
finite-program theorem. Repair it to $\widetilde z$ at $m'$ as in the
lemma. Since $L(m')$ is a minimum over all obedient arrays at $m'$,
\begin{align*}
 L(m')&\leq\sum w^{m'}\widetilde z\\
 &\leq\sum w^{m'}z+2W\alpha\\
 &\leq\sum w^mz+L_wh+2W\alpha
   =L(m)+\omega(h).
\end{align*}
Sums are over all state-action pairs. Interchanging $m$ and $m'$ proves
the absolute bound. The displayed linear bound uses
$\sigma+2L_uh\geq\sigma$. Thus $L$ is continuous; a continuous real
function on the nonempty compact set $\mathcal M$ attains its maximum.
\end{proof}

\begin{theorem}[A verifiable finite-net design guarantee]
Let $N_h\subseteq\mathcal M$ be a finite $h$-net: every member of
$\mathcal M$ is at distance at most $h$ from some $m\in N_h$.
For each net point suppose the primal and dual certificates give
\[
 \ell_m\leq L(m)\leq u_m,\qquad u_m-\ell_m\leq\varepsilon_0.
\]
Select $\widehat m\in N_h$ with largest $\ell_m$. Then
\[
 L(\widehat m)\geq\max_{m\in\mathcal M}L(m)
                       -\omega(h)-\varepsilon_0.
\]
A certified $h$-net, payoff Lipschitz bounds, a verified uniform margin
$\sigma$, and the finite primal/dual inequalities together give a
verifiable guarantee. Compactness establishes the existence of finite
nets; computing and certifying a net from an unspecified abstract
parameter space requires additional input.
\end{theorem}
\begin{proof}
Let $m^*$ maximize $L$ and choose a net point $m_0$ within distance
$h$. The continuity bound gives $L(m_0)\geq L(m^*)-\omega(h)$.
The finite-library certificate argument yields
\[
 L(\widehat m)\geq\ell_{\widehat m}\geq\ell_{m_0}
     \geq L(m_0)-\varepsilon_0.
\]
Combining the inequalities proves the result. Every net point lies in
$\mathcal M$, so its feasibility is part of the declared parameter
class rather than inferred from its proximity to another mechanism.
\end{proof}

For a requested $\varepsilon>0$, write
$C=L_w+4WL_u/\sigma$. If $C>0$, choose
$h\leq\varepsilon/(2C)$ and $\varepsilon_0\leq\varepsilon/2$.
The guarantee is then $\varepsilon$-optimal. If $C=0$, the continuity
bound makes $L$ constant on the class, so any member suffices. With
rational coefficients at rational net points, the primal and dual
certificates can be checked by exact arithmetic as in Version 1.
This does not assert a polynomial bound on net size or linear-program
bit complexity in an arbitrary mechanism dimension.

\section{The margin is nonvacuous and cannot be silently omitted}
\begin{example}[A continuous class with a common strict certificate]
Take one decision-maker, two equiprobable states $\omega\in\{0,1\}$,
actions $a\in\{0,1\}$, and parameter $m\in[-1/2,1/2]$. Let
\[
 u^m(a,\omega)=\mathbf1_{\{a=\omega\}}+m\mathbf1_{\{a=1\}},
 \qquad w^m(a,\omega)=-m\mathbf1_{\{a=1\}}.
\]
One interpretation is a bounded action-dependent subsidy, where the
underlying state-dependent return is the first utility term. Recommend
the true state: $\bar z(0,0)=\bar z(1,1)=1/2$ and the other masses
are zero. The obedience slack for recommendation zero is
$(1-m)/2\geq1/4$; for recommendation one it is
$(1+m)/2\geq1/4$. Thus $\sigma=1/4$ works for the whole interval,
with $L_u=L_w=1$ and $W=1/2$. The example demonstrates a uniform strict
certificate for a class with changing agent incentives. Any additional
participation or budget requirement must still be specified; this
example imposes only the finite-action and bounded-transfer conventions.
\end{example}

\begin{proposition}[A jump at an indifference point]
Without uniform strict feasibility, even continuous bounded payoff
arrays on a compact parameter interval need not produce a continuous
worst-equilibrium value. With one state, one player, actions $0,1$,
parameter $m\in[0,1]$, agent utilities $u^m(0)=0$, $u^m(1)=m$,
and designer payoffs $w^m(0)=0$, $w^m(1)=1$, one has
\[
 L(0)=0,\qquad L(m)=1\quad(m>0).
\]
\end{proposition}
\begin{proof}
There is no payoff-state uncertainty to reveal. At $m>0$, action one
strictly dominates action zero. Its obedience constraint against
recommending zero forces the mass on action zero to be zero, so every
obedient outcome gives designer payoff one. At $m=0$, either action
is optimal; recommending zero with probability one is obedient and
gives payoff zero, which is the minimum. These values establish the
jump despite the Lipschitz agent and designer arrays.
\end{proof}
This counterexample does not show that every particular dense grid
fails. It shows why continuity or a uniform approximation error cannot
be obtained merely by assuming continuous primitives and compact
parameters. The repair proof states an explicit condition under which
the needed quantitative stability is valid. It is an application of
standard finite linear-program stability reasoning, not a claim that
uniform strict obedience holds in arbitrary mechanism classes.

\section{Remaining gap}
The reduction includes pessimistic equilibrium selection without presuming truthful reporting. If the catalogue ambiguity set restricts signal alphabets or preserves pre-existing private information, the unrestricted program can be only a conservative bound unless those restrictions are encoded. Version 2 supplies an outer approximation theorem for a compact class with fixed message sizes and a uniform strict-obedience certificate. It does not approximate arbitrary message alphabets, prove that every optimal mechanism has a strict certificate, or automatically preserve binding interim robust constraints under discretization. Enlarging any ambiguity set weakly lowers the value of each fixed mechanism; when robust feasibility is also tightened, taking the supremum over the smaller feasible class preserves that inequality.

\section{Completion Estimate}
45\%

This is a heuristic estimate for the full catalogue target.
The continuous-class theorem now combines exact pessimistic-equilibrium certificates with a quantitative finite-net approximation and an explicit failure example. Uniform strict obedience, fixed message sets, feasible net construction and the unrestricted common-prior signal ambiguity remain substantial restrictions on the full design target.

\begin{thebibliography}{9}
\bibitem{BM} D. Bergemann and S. Morris (2016), Bayes correlated equilibrium and the comparison of information structures in games, \emph{Theoretical Economics} 11, 487--522. \url{https://www.econtheory.org/ojs/index.php/te/article/viewArticle/20160487}. Source of the obedience/outcome characterization; the finite-library certification formulation above is derived here.
\bibitem{robust} D. Bergemann and S. Morris (2005), Robust mechanism design, \emph{Econometrica} 73, 1771--1813. \url{https://doi.org/10.1111/j.1468-0262.2005.00638.x}. Background on informational robustness, not a solution of arbitrary outer design problems.
\end{thebibliography}
\end{document}
